\section{Introduction}\label{sec:intro}

Indian agricultural wholesale markets (\emph{mandis}) report minimum, maximum and modal prices by commodity and variety through AGMARKNET and public data mirrors.
These reports provide price information for farmers comparing markets and sale timing, but their value as forecasting data depends on reporting continuity, source completeness and the evaluation protocol.
The empirical question is how much a learned forecast improves on simply retaining the latest reported price, and how cautiously that improvement should be interpreted.

The mandi forecasting literature includes studies of individual markets, small groups of markets and monthly indices (Section~\ref{sec:related}), often comparing deep or hybrid models with ARIMA \cite{paul2022brinjal,manogna2025deep,kumar2021onion}.
Forecast-evaluation research motivates a complementary emphasis: strong simple benchmarks \cite{makridakis2018statistical,makridakis2020m4}, appropriate scaling across heterogeneous series \cite{hyndman2006another}, alignment between the training loss and the scoring function \cite{gneiting2011making,kolassa2016evaluating}, and explicit treatment of leakage and data vintage \cite{hewamalage2022pitfalls,bokun2023fredsd}.
These considerations motivate an auditable empirical study rather than a claim of model superiority across published studies with different data and protocols.

Our primary contribution is a smaller, resource-limited \emph{clean-source release study} of daily mandi modal-price forecasting.
For each horizon and filtering regime, the learner uses \num{500000} deterministically sampled training rows from 2024; January--June 2025 is used for model selection and early stopping, July--December 2025 for interval calibration, and January--October 2026 for retrospective evaluation.
Targets are purged at period boundaries.
The release documents the feature builds and saved predictions, but retains a fixed historical panel as an external input.
It therefore supports claims about this release and observed panel, rather than validation of the earlier full national benchmark.

We make three empirical contributions, with their limits stated alongside the findings.
\begin{enumerate}
  \item \textbf{A chronological, baseline-explicit release comparison.}
    We compare an absolute-loss LightGBM candidate with persistence and four other simple baselines at the next-report, 3-, 7- and 14-day horizons.
    The two regimes exclude or retain reports flagged by the fivefold jump rule while preserving the other validity checks and panel membership.
    Across the eight fits, pooled MAE is 4.4--5.9\% lower than persistence on the retrospective evaluation rows.
    These are descriptive historical gains; adoption of LightGBM remains exploratory, with no independent confirmation of its superiority over the earlier Spark models or of benefits to farmer decisions.
  \item \textbf{Explicit separation of pooled error and scaled error.}
    Pooled MAE and MASE answer different questions: MASE divides errors by each series' training-period, horizon-specific naive error scale, whereas the pooled MAE comparison uses evaluation-period persistence error as its reference.
    MASE also omits approximately 4.4--4.5\% of forecast rows because their training scales are not positive and finite.
    It consequently describes a different population from full-row MAE; a MASE value above one is compatible with a pooled MAE gain over evaluation-period persistence.
  \item \textbf{Dependence-aware sensitivity analysis and interval assessment.}
    We report calendar-day block-resampling and dependence sensitivities alongside retrospective prediction-interval coverage.
    Forecast rows are not independent replicates, and the inference relies on approximate weak-dependence and stationarity assumptions conditional on the fitted models and observed panel.
    Neither small test probabilities nor narrow resampling intervals establish validated inferential coverage or generalisation to independent new markets.
\end{enumerate}

An earlier source-completeness audit provides historical supporting evidence that incomplete public mirrors can alter apparent forecast accuracy and interval coverage.
That audit is not reproduced as the principal validation of this release, and the release does not establish the original full benchmark's source state.
Similarly, the 2026 evaluation is retrospective: its forecasts were not issued as the reports arrived, so it does not constitute prospective confirmation.
Independent confirmation, expert statistical review and future-issued evaluation remain outstanding.
